\documentclass[12pt, a4paper]{article}

% Packages
\usepackage[utf8]{inputenc}
\usepackage[T1]{fontenc}
\usepackage[english]{babel}
\usepackage{amsmath, amssymb, amsthm, amsfonts}
\usepackage{geometry}
\usepackage{hyperref}
\usepackage{bookmark}
\usepackage{enumitem}
\usepackage{fancyhdr}
\usepackage{titlesec}
\usepackage{xcolor}

% Geometry setup
\geometry{
    top=2.5cm,
    bottom=2.5cm,
    left=2.5cm,
    right=2.5cm
}
\setlength{\headheight}{30pt}

% Hyperref setup
\hypersetup{
    colorlinks=true,
    linkcolor=blue,
    filecolor=magenta,      
    urlcolor=cyan,
}

% Theorem environments
\newtheorem{theorem}{Theorem}
\newtheorem{proposition}[theorem]{Proposition}
\newtheorem{lemma}[theorem]{Lemma}
\newtheorem{corollary}[theorem]{Corollary}

\theoremstyle{definition}
\newtheorem{definition}{Definition}
\newtheorem{exercise}{Exercise}
\newtheorem{remark}{Remark}

% Layout/Style
\pagestyle{fancy}
\fancyhf{}
\rhead{Stochastic Finance - TDs 1 \& 2}
\lhead{Corrections \& Notes}
\cfoot{\thepage}

\title{\textbf{Stochastic Finance and Risk Modelling} \\ \large TDs 1 \& 2: Exercises and Corrections}
\author{Pedro}
\date{\today}

\begin{document}

\maketitle

\tableofcontents

\vspace{1cm}

\section*{Introduction}
\addcontentsline{toc}{section}{Introduction}
This document contains the step-by-step corrections, notes, and detailed explanations for the exercises in \textbf{TD 1 and TD 2} of the Stochastic Finance and Risk Modelling course.

\section{TD 1}

\begin{exercise}[Conditional expectation]
Let $X$ be a r.v. taking values in some finite set and $Y \in L^1(\Omega, \mathcal{F}, \mathbb{P})$. Prove that there exist $n \in \mathbb{N}$ and a partition $(A_1, \dots, A_n)$ of $\Omega$ s.t.
\[
\mathbb{E}[Y|X] = \sum_{i=1}^n \frac{\mathbb{E}[Y \mathbf{1}_{A_i}]}{\mathbb{P}[A_i]}\mathbf{1}_{A_i} \quad \mathbb{P}\text{-a.s.}
\]
\end{exercise}

\textbf{Correction \& Explanation:}

\textbf{Step 1: Construct the Partition and the Sub-$\sigma$-algebra.} \\
Since $X$ assumes a finite set of distinct values $\{x_1, x_2, \dots, x_n\}$, we can define the event $A_i$ as the set of outcomes where $X$ takes the value $x_i$:
\[ A_i = \{\omega \in \Omega : X(\omega) = x_i\} \]
Because every outcome $\omega$ maps to exactly one value $X(\omega)$, the sets $A_1, \dots, A_n$ are disjoint and their union is the entire sample space $\Omega$. Therefore, $(A_1, \dots, A_n)$ is a \textbf{partition} of $\Omega$.
The $\sigma$-algebra generated by $X$, denoted by $\mathcal{G} = \sigma(X)$, is simply the collection of all possible unions of these blocks $A_i$.

\textbf{Step 2: Understand the Natural Form of $\mathbb{E}[Y|X]$.} \\
By definition, $Z = \mathbb{E}[Y|X]$ must be $\sigma(X)$-measurable. This means that $Z$ can only change its value when $X$ changes its value. Since $X$ is constant on each block $A_i$, the variable $Z$ must also be \textbf{constant within each block $A_i$}. 
Thus, $Z$ must take the form of a simple function:
\[ Z = \sum_{i=1}^n c_i \mathbf{1}_{A_i} \]
where $c_i$ are constants that we need to determine.

\textbf{Step 3: Use the Partial Averaging Property.} \\
The fundamental property of the conditional expectation states that for any event $A \in \sigma(X)$, we must have:
\[ \mathbb{E}[Z \mathbf{1}_{A}] = \mathbb{E}[Y \mathbf{1}_{A}] \]
In particular, this must hold for each block $A_i$ in our partition. Let's calculate the left side by substituting our expression for $Z$:
\[ \mathbb{E}[Z \mathbf{1}_{A_i}] = \mathbb{E}\left[\left( \sum_{j=1}^n c_j \mathbf{1}_{A_j} \right) \mathbf{1}_{A_i}\right] \]
Since the sets $A_i$ and $A_j$ are disjoint for $i \neq j$, the product of their indicator functions is zero ($\mathbf{1}_{A_j}\mathbf{1}_{A_i} = 0$). The only surviving term is when $j = i$, where $\mathbf{1}_{A_i} \mathbf{1}_{A_i} = \mathbf{1}_{A_i}$:
\[ \mathbb{E}[Z \mathbf{1}_{A_i}] = \mathbb{E}[c_i \mathbf{1}_{A_i}] = c_i \mathbb{E}[\mathbf{1}_{A_i}] = c_i \mathbb{P}(A_i) \]

Now, we equate this back to the initial property:
\[ c_i \mathbb{P}(A_i) = \mathbb{E}[Y \mathbf{1}_{A_i}] \]
Assuming $\mathbb{P}(A_i) > 0$, we can isolate $c_i$:
\[ c_i = \frac{\mathbb{E}[Y \mathbf{1}_{A_i}]}{\mathbb{P}(A_i)} \]
(If $\mathbb{P}(A_i) = 0$, the value of $c_i$ does not affect the almost sure equality).

\textbf{Conclusion} \\
Substituting the constants $c_i$ back into the generic form of $Z$, we obtain:
\[ \mathbb{E}[Y|X] = Z = \sum_{i=1}^n \frac{\mathbb{E}[Y \mathbf{1}_{A_i}]}{\mathbb{P}(A_i)} \mathbf{1}_{A_i} \quad \mathbb{P}\text{-a.s.} \]
This completes the proof.



\end{document}
